% --- Archon physics formalization source begin ---
% archon:physics
% archon:covers PhyXMiniProblems/problem_phyx_mini_0618.lean
% archon:source-report reports/phyx_mini/problem_phyx_mini_0618.source.json

\chapter{Physics problem phyx\_mini\_0618}
\label{ch:PhyXMiniProblems_problem_phyx_mini_0618}

\paragraph{Problem source.}
\#\# Physical scenario

\textbackslash{}(\{\}\textasciicircum{}\{232\}\_\{92\}\textbackslash{}mathrm\{U\}\textbackslash{}) (mass \textbackslash{}(=232.037156\textbackslash{},\textbackslash{}mathrm\{u\}\textbackslash{})) decays to \textbackslash{}(\{\}\textasciicircum{}\{228\}\_\{90\}\textbackslash{}mathrm\{Th\}\textbackslash{}) (\textbackslash{}(228.028741\textbackslash{},\textbackslash{}mathrm\{u\}\textbackslash{})) with the emission of an \textbackslash{}(\textbackslash{}alpha\textbackslash{}) particle. (As always, masses given are for neutral atoms.)

\#\# Question

Calculate the disintegration energy.

\#\# Answer choices

- A: A: 5.0MeV

- B: B: 5.2MeV

- C: C: 5.4MeV

- D: D: 4.8MeV

\#\# Figure caption (auxiliary; use the image as primary evidence)

The image shows a physics diagram illustrating the motion of an alpha particle and a daughter nucleus, likely depicting a nuclear reaction. There are two objects represented as small pink circles. 

- The object on the left is labeled "α particle" with its velocity represented by a leftward green arrow labeled \textbackslash{}( \textbackslash{}vec\{V\}\_α \textbackslash{}), and its mass indicated as \textbackslash{}( m\_α \textbackslash{}).
- The object on the right is labeled "Daughter nucleus" with its velocity represented by a rightward green arrow labeled \textbackslash{}( \textbackslash{}vec\{V\}\_D \textbackslash{}), and its mass indicated as \textbackslash{}( m\_D \textbackslash{}).

The arrows show the direction of motion for each particle following a likely nuclear decay event.

\#\# Dataset metadata

Nuclear Physics; Physical Model Grounding Reasoning, Multi-Formula Reasoning

\paragraph{Recorded answer/context.}
C: C: 5.4MeV

\paragraph{Figure/image path.}
/root/proposal\_for\_physic/science-mango/phyx\_mini\_run/phyx\_data/test\_image/618.png

\paragraph{Formalization target.}
create a compiling Lean file with sorry bodies at `PhyXMiniProblems/problem\_phyx\_mini\_0618.lean`. The Lean declarations must preserve the physical quantities, dimensions or dimensional roles, figure labels, governing-law hypotheses, and final relation expressed by this problem.
Use Mathlib/Physlib names found through LeanExplore where available. If a physics API is missing, introduce faithful local abstractions rather than scalar placeholder aliases.

\begin{theorem}[Physics formalization target]
\label{thm:physics:phyx_mini_0618:target}
\lean{PhyXMiniProblems.ProblemPhyXMini0618.problem_phyx_mini_0618}
\uses{def:physics:phyx-mini-0618:phyxminiproblems-problemphyxmini0618-atomicmassunitinkilograms, def:physics:phyx-mini-0618:phyxminiproblems-problemphyxmini0618-massinatomicmassunits, def:physics:phyx-mini-0618:phyxminiproblems-problemphyxmini0618-electronvoltinjoules, def:physics:phyx-mini-0618:phyxminiproblems-problemphyxmini0618-energyinmegaelectronvolts, def:physics:phyx-mini-0618:phyxminiproblems-problemphyxmini0618-speedoflightinmeterspersecond, def:physics:phyx-mini-0618:phyxminiproblems-problemphyxmini0618-uranium232alphadecaysetup, def:physics:phyx-mini-0618:phyxminiproblems-problemphyxmini0618-matchesalphadecayproblemdata, def:physics:phyx-mini-0618:phyxminiproblems-problemphyxmini0618-matchesreferenceatomicdata, def:physics:phyx-mini-0618:phyxminiproblems-problemphyxmini0618-matchesprimaryalphadecayfigure, def:physics:phyx-mini-0618:phyxminiproblems-problemphyxmini0618-hasphysicalalphadecayparameters, def:physics:phyx-mini-0618:phyxminiproblems-problemphyxmini0618-satisfiesalphadecayenergylaws, def:physics:phyx-mini-0618:phyxminiproblems-problemphyxmini0618-recordedanswerchoice, def:physics:phyx-mini-0618:phyxminiproblems-problemphyxmini0618-matchesdisplayeddisintegrationenergy}
The assigned autoformalize agent should translate this physics problem into Lean declarations in the covered file, with theorem and lemma proof bodies written as `by sorry`.
\end{theorem}
\begin{proof}
This is an autoformalization task, not a proof task. Produce statements that can later be proved by the physics prover without weakening the physical model.
\end{proof}

% --- Archon declaration topology begin ---
\section{Declaration topology}
\begin{definition}[Mass Quantity]
  \label{def:physics:phyx-mini-0618:phyxminiproblems-problemphyxmini0618-massquantity}
  \lean{PhyXMiniProblems.ProblemPhyXMini0618.MassQuantity}
  A nonnegative, unit-independent physical rest mass
\end{definition}
\begin{proof}
  This is the declared model definition of Mass Quantity.
\end{proof}
\begin{definition}[Planar Velocity Quantity]
  \label{def:physics:phyx-mini-0618:phyxminiproblems-problemphyxmini0618-planarvelocityquantity}
  \lean{PhyXMiniProblems.ProblemPhyXMini0618.PlanarVelocityQuantity}
  A unit-independent velocity vector in the two-dimensional plane of the supplied recoil diagram
\end{definition}
\begin{proof}
  This is the declared model definition of Planar Velocity Quantity.
\end{proof}
\begin{definition}[Diagram Axis]
  \label{def:physics:phyx-mini-0618:phyxminiproblems-problemphyxmini0618-diagramaxis}
  \lean{PhyXMiniProblems.ProblemPhyXMini0618.DiagramAxis}
  The two coordinate directions in the recoil diagram
\end{definition}
\begin{proof}
  This is the declared model definition of Diagram Axis.
\end{proof}
\begin{definition}[to Fin]
  \label{def:physics:phyx-mini-0618:phyxminiproblems-problemphyxmini0618-diagramaxis-tofin}
  \lean{PhyXMiniProblems.ProblemPhyXMini0618.DiagramAxis.toFin}
  \uses{def:physics:phyx-mini-0618:phyxminiproblems-problemphyxmini0618-diagramaxis}
  Coordinate index associated with a diagram axis
\end{definition}
\begin{proof}
  This is the declared model definition of to Fin.
\end{proof}
\begin{definition}[mass In Kilograms]
  \label{def:physics:phyx-mini-0618:phyxminiproblems-problemphyxmini0618-massinkilograms}
  \lean{PhyXMiniProblems.ProblemPhyXMini0618.massInKilograms}
  \uses{def:physics:phyx-mini-0618:phyxminiproblems-problemphyxmini0618-massquantity}
  Kilogram readout of a physical mass in coherent SI units
\end{definition}
\begin{proof}
  This is the declared model definition of mass In Kilograms.
\end{proof}
\begin{definition}[atomic Mass Unit In Kilograms]
  \label{def:physics:phyx-mini-0618:phyxminiproblems-problemphyxmini0618-atomicmassunitinkilograms}
  \lean{PhyXMiniProblems.ProblemPhyXMini0618.atomicMassUnitInKilograms}
  The SI kilogram value used for one unified atomic mass unit
\end{definition}
\begin{proof}
  This is the declared model definition of atomic Mass Unit In Kilograms.
\end{proof}
\begin{definition}[mass In Atomic Mass Units]
  \label{def:physics:phyx-mini-0618:phyxminiproblems-problemphyxmini0618-massinatomicmassunits}
  \lean{PhyXMiniProblems.ProblemPhyXMini0618.massInAtomicMassUnits}
  \uses{def:physics:phyx-mini-0618:phyxminiproblems-problemphyxmini0618-massquantity, def:physics:phyx-mini-0618:phyxminiproblems-problemphyxmini0618-massinkilograms, def:physics:phyx-mini-0618:phyxminiproblems-problemphyxmini0618-atomicmassunitinkilograms}
  Unified-atomic-mass-unit readout of a physical mass
\end{definition}
\begin{proof}
  This is the declared model definition of mass In Atomic Mass Units.
\end{proof}
\begin{definition}[velocity Component In Meters Per Second]
  \label{def:physics:phyx-mini-0618:phyxminiproblems-problemphyxmini0618-velocitycomponentinmeterspersecond}
  \lean{PhyXMiniProblems.ProblemPhyXMini0618.velocityComponentInMetersPerSecond}
  \uses{def:physics:phyx-mini-0618:phyxminiproblems-problemphyxmini0618-planarvelocityquantity, def:physics:phyx-mini-0618:phyxminiproblems-problemphyxmini0618-diagramaxis}
  One component of a planar velocity, read in metres per second
\end{definition}
\begin{proof}
  This is the declared model definition of velocity Component In Meters Per Second.
\end{proof}
\begin{definition}[energy In Joules]
  \label{def:physics:phyx-mini-0618:phyxminiproblems-problemphyxmini0618-energyinjoules}
  \lean{PhyXMiniProblems.ProblemPhyXMini0618.energyInJoules}
  Joule readout of a dimensionful energy
\end{definition}
\begin{proof}
  This is the declared model definition of energy In Joules.
\end{proof}
\begin{definition}[electron Volt In Joules]
  \label{def:physics:phyx-mini-0618:phyxminiproblems-problemphyxmini0618-electronvoltinjoules}
  \lean{PhyXMiniProblems.ProblemPhyXMini0618.electronVoltInJoules}
  Physlib's exact electron volt, read in joules
\end{definition}
\begin{proof}
  This is the declared model definition of electron Volt In Joules.
\end{proof}
\begin{definition}[energy In Mega Electron Volts]
  \label{def:physics:phyx-mini-0618:phyxminiproblems-problemphyxmini0618-energyinmegaelectronvolts}
  \lean{PhyXMiniProblems.ProblemPhyXMini0618.energyInMegaElectronVolts}
  \uses{def:physics:phyx-mini-0618:phyxminiproblems-problemphyxmini0618-energyinjoules, def:physics:phyx-mini-0618:phyxminiproblems-problemphyxmini0618-electronvoltinjoules}
  A physical energy expressed as a numerical number of megaelectronvolts
\end{definition}
\begin{proof}
  This is the declared model definition of energy In Mega Electron Volts.
\end{proof}
\begin{definition}[speed Of Light In Meters Per Second]
  \label{def:physics:phyx-mini-0618:phyxminiproblems-problemphyxmini0618-speedoflightinmeterspersecond}
  \lean{PhyXMiniProblems.ProblemPhyXMini0618.speedOfLightInMetersPerSecond}
  Physlib's exact vacuum speed of light, read in metres per second
\end{definition}
\begin{proof}
  This is the declared model definition of speed Of Light In Meters Per Second.
\end{proof}
\begin{definition}[Nuclide]
  \label{def:physics:phyx-mini-0618:phyxminiproblems-problemphyxmini0618-nuclide}
  \lean{PhyXMiniProblems.ProblemPhyXMini0618.Nuclide}
  The three nuclides used in the decay and neutral-atom mass balance
\end{definition}
\begin{proof}
  This is the declared model definition of Nuclide.
\end{proof}
\begin{definition}[atomic Number]
  \label{def:physics:phyx-mini-0618:phyxminiproblems-problemphyxmini0618-nuclide-atomicnumber}
  \lean{PhyXMiniProblems.ProblemPhyXMini0618.Nuclide.atomicNumber}
  \uses{def:physics:phyx-mini-0618:phyxminiproblems-problemphyxmini0618-nuclide}
  Proton number Z of each nuclide
\end{definition}
\begin{proof}
  This is the declared model definition of atomic Number.
\end{proof}
\begin{definition}[mass Number]
  \label{def:physics:phyx-mini-0618:phyxminiproblems-problemphyxmini0618-nuclide-massnumber}
  \lean{PhyXMiniProblems.ProblemPhyXMini0618.Nuclide.massNumber}
  \uses{def:physics:phyx-mini-0618:phyxminiproblems-problemphyxmini0618-nuclide}
  Nucleon (mass) number A of each nuclide
\end{definition}
\begin{proof}
  This is the declared model definition of mass Number.
\end{proof}
\begin{definition}[Alpha Decay Channel]
  \label{def:physics:phyx-mini-0618:phyxminiproblems-problemphyxmini0618-alphadecaychannel}
  \lean{PhyXMiniProblems.ProblemPhyXMini0618.AlphaDecayChannel}
  \uses{def:physics:phyx-mini-0618:phyxminiproblems-problemphyxmini0618-nuclide}
  The parent, daughter, and emitted nuclide identities in one nuclear decay channel. The emitted helium-4 nucleus is the physical alpha particle
\end{definition}
\begin{proof}
  This is the declared model definition of Alpha Decay Channel.
\end{proof}
\begin{definition}[Conserves Nuclear Numbers]
  \label{def:physics:phyx-mini-0618:phyxminiproblems-problemphyxmini0618-conservesnuclearnumbers}
  \lean{PhyXMiniProblems.ProblemPhyXMini0618.ConservesNuclearNumbers}
  \uses{def:physics:phyx-mini-0618:phyxminiproblems-problemphyxmini0618-alphadecaychannel}
  Conservation of nucleon number and proton number in a decay channel
\end{definition}
\begin{proof}
  This is the declared model definition of Conserves Nuclear Numbers.
\end{proof}
\begin{definition}[Decay Product]
  \label{def:physics:phyx-mini-0618:phyxminiproblems-problemphyxmini0618-decayproduct}
  \lean{PhyXMiniProblems.ProblemPhyXMini0618.DecayProduct}
  The two physical products represented by circles in the supplied image
\end{definition}
\begin{proof}
  This is the declared model definition of Decay Product.
\end{proof}
\begin{definition}[Product Figure Label]
  \label{def:physics:phyx-mini-0618:phyxminiproblems-problemphyxmini0618-productfigurelabel}
  \lean{PhyXMiniProblems.ProblemPhyXMini0618.ProductFigureLabel}
  Textual object labels visible above the two product markers
\end{definition}
\begin{proof}
  This is the declared model definition of Product Figure Label.
\end{proof}
\begin{definition}[Mass Figure Label]
  \label{def:physics:phyx-mini-0618:phyxminiproblems-problemphyxmini0618-massfigurelabel}
  \lean{PhyXMiniProblems.ProblemPhyXMini0618.MassFigureLabel}
  Symbolic mass labels printed below the two products
\end{definition}
\begin{proof}
  This is the declared model definition of Mass Figure Label.
\end{proof}
\begin{definition}[Velocity Figure Label]
  \label{def:physics:phyx-mini-0618:phyxminiproblems-problemphyxmini0618-velocityfigurelabel}
  \lean{PhyXMiniProblems.ProblemPhyXMini0618.VelocityFigureLabel}
  Symbolic velocity-vector labels printed beside the arrows
\end{definition}
\begin{proof}
  This is the declared model definition of Velocity Figure Label.
\end{proof}
\begin{definition}[Horizontal Direction]
  \label{def:physics:phyx-mini-0618:phyxminiproblems-problemphyxmini0618-horizontaldirection}
  \lean{PhyXMiniProblems.ProblemPhyXMini0618.HorizontalDirection}
  Qualitative horizontal arrow directions in the supplied image
\end{definition}
\begin{proof}
  This is the declared model definition of Horizontal Direction.
\end{proof}
\begin{definition}[Marker Color]
  \label{def:physics:phyx-mini-0618:phyxminiproblems-problemphyxmini0618-markercolor}
  \lean{PhyXMiniProblems.ProblemPhyXMini0618.MarkerColor}
  Marker colors used in the raster
\end{definition}
\begin{proof}
  This is the declared model definition of Marker Color.
\end{proof}
\begin{definition}[Arrow Color]
  \label{def:physics:phyx-mini-0618:phyxminiproblems-problemphyxmini0618-arrowcolor}
  \lean{PhyXMiniProblems.ProblemPhyXMini0618.ArrowColor}
  Velocity-arrow colors used in the raster
\end{definition}
\begin{proof}
  This is the declared model definition of Arrow Color.
\end{proof}
\begin{definition}[Alpha Decay Figure]
  \label{def:physics:phyx-mini-0618:phyxminiproblems-problemphyxmini0618-alphadecayfigure}
  \lean{PhyXMiniProblems.ProblemPhyXMini0618.AlphaDecayFigure}
  \uses{def:physics:phyx-mini-0618:phyxminiproblems-problemphyxmini0618-decayproduct, def:physics:phyx-mini-0618:phyxminiproblems-problemphyxmini0618-productfigurelabel, def:physics:phyx-mini-0618:phyxminiproblems-problemphyxmini0618-massfigurelabel, def:physics:phyx-mini-0618:phyxminiproblems-problemphyxmini0618-velocityfigurelabel, def:physics:phyx-mini-0618:phyxminiproblems-problemphyxmini0618-horizontaldirection, def:physics:phyx-mini-0618:phyxminiproblems-problemphyxmini0618-markercolor, def:physics:phyx-mini-0618:phyxminiproblems-problemphyxmini0618-arrowcolor}
  Presentation-level evidence transcribed from image 618. Its horizontal coordinates preserve only left-to-right ordering and are not physical distances
\end{definition}
\begin{proof}
  This is the declared model definition of Alpha Decay Figure.
\end{proof}
\begin{definition}[Uranium232 Alpha Decay Setup]
  \label{def:physics:phyx-mini-0618:phyxminiproblems-problemphyxmini0618-uranium232alphadecaysetup}
  \lean{PhyXMiniProblems.ProblemPhyXMini0618.Uranium232AlphaDecaySetup}
  \uses{def:physics:phyx-mini-0618:phyxminiproblems-problemphyxmini0618-massquantity, def:physics:phyx-mini-0618:phyxminiproblems-problemphyxmini0618-planarvelocityquantity, def:physics:phyx-mini-0618:phyxminiproblems-problemphyxmini0618-nuclide, def:physics:phyx-mini-0618:phyxminiproblems-problemphyxmini0618-alphadecaychannel, def:physics:phyx-mini-0618:phyxminiproblems-problemphyxmini0618-decayproduct, def:physics:phyx-mini-0618:phyxminiproblems-problemphyxmini0618-alphadecayfigure}
  All independent physical quantities in the decay setup. The neutral-atom masses used for Q-value bookkeeping are kept distinct from the actual bare alpha-particle and daughter-nucleus masses labelled in the figure. The mass defect and disintegration energy are independent fields; neither is defined from a displayed answer
\end{definition}
\begin{proof}
  This is the declared model definition of Uranium232 Alpha Decay Setup.
\end{proof}
\begin{definition}[product Rest Mass]
  \label{def:physics:phyx-mini-0618:phyxminiproblems-problemphyxmini0618-productrestmass}
  \lean{PhyXMiniProblems.ProblemPhyXMini0618.productRestMass}
  \uses{def:physics:phyx-mini-0618:phyxminiproblems-problemphyxmini0618-massquantity, def:physics:phyx-mini-0618:phyxminiproblems-problemphyxmini0618-decayproduct, def:physics:phyx-mini-0618:phyxminiproblems-problemphyxmini0618-uranium232alphadecaysetup}
  Physical rest mass associated with a product's m\_alpha or m\_D label
\end{definition}
\begin{proof}
  This is the declared model definition of product Rest Mass.
\end{proof}
\begin{definition}[Matches Alpha Decay Problem Data]
  \label{def:physics:phyx-mini-0618:phyxminiproblems-problemphyxmini0618-matchesalphadecayproblemdata}
  \lean{PhyXMiniProblems.ProblemPhyXMini0618.MatchesAlphaDecayProblemData}
  \uses{def:physics:phyx-mini-0618:phyxminiproblems-problemphyxmini0618-massinatomicmassunits, def:physics:phyx-mini-0618:phyxminiproblems-problemphyxmini0618-uranium232alphadecaysetup}
  The decay identities and the two neutral-atom masses explicitly supplied in the prose. Decimal values are represented exactly by rational numbers
\end{definition}
\begin{proof}
  This is the declared model definition of Matches Alpha Decay Problem Data.
\end{proof}
\begin{definition}[Matches Reference Atomic Data]
  \label{def:physics:phyx-mini-0618:phyxminiproblems-problemphyxmini0618-matchesreferenceatomicdata}
  \lean{PhyXMiniProblems.ProblemPhyXMini0618.MatchesReferenceAtomicData}
  \uses{def:physics:phyx-mini-0618:phyxminiproblems-problemphyxmini0618-massinatomicmassunits, def:physics:phyx-mini-0618:phyxminiproblems-problemphyxmini0618-uranium232alphadecaysetup}
  Calibrated helium-4 neutral-atom mass required by the neutral-atom Q-value calculation. This is reference nuclear data, not the requested energy
\end{definition}
\begin{proof}
  This is the declared model definition of Matches Reference Atomic Data.
\end{proof}
\begin{definition}[Matches Primary Alpha Decay Figure]
  \label{def:physics:phyx-mini-0618:phyxminiproblems-problemphyxmini0618-matchesprimaryalphadecayfigure}
  \lean{PhyXMiniProblems.ProblemPhyXMini0618.MatchesPrimaryAlphaDecayFigure}
  \uses{def:physics:phyx-mini-0618:phyxminiproblems-problemphyxmini0618-velocitycomponentinmeterspersecond, def:physics:phyx-mini-0618:phyxminiproblems-problemphyxmini0618-uranium232alphadecaysetup}
  Labels, marker colors, ordering, and outward velocity arrows read from the primary bitmap. The component signs connect the drawn vector labels to the dimensionful velocities without assigning their magnitudes
\end{definition}
\begin{proof}
  This is the declared model definition of Matches Primary Alpha Decay Figure.
\end{proof}
\begin{definition}[Has Physical Alpha Decay Parameters]
  \label{def:physics:phyx-mini-0618:phyxminiproblems-problemphyxmini0618-hasphysicalalphadecayparameters}
  \lean{PhyXMiniProblems.ProblemPhyXMini0618.HasPhysicalAlphaDecayParameters}
  \uses{def:physics:phyx-mini-0618:phyxminiproblems-problemphyxmini0618-massinkilograms, def:physics:phyx-mini-0618:phyxminiproblems-problemphyxmini0618-atomicmassunitinkilograms, def:physics:phyx-mini-0618:phyxminiproblems-problemphyxmini0618-energyinjoules, def:physics:phyx-mini-0618:phyxminiproblems-problemphyxmini0618-electronvoltinjoules, def:physics:phyx-mini-0618:phyxminiproblems-problemphyxmini0618-speedoflightinmeterspersecond, def:physics:phyx-mini-0618:phyxminiproblems-problemphyxmini0618-uranium232alphadecaysetup, def:physics:phyx-mini-0618:phyxminiproblems-problemphyxmini0618-productrestmass}
  Positivity and nondegeneracy conditions for the physical decay branch
\end{definition}
\begin{proof}
  This is the declared model definition of Has Physical Alpha Decay Parameters.
\end{proof}
\begin{definition}[Satisfies Alpha Decay Energy Laws]
  \label{def:physics:phyx-mini-0618:phyxminiproblems-problemphyxmini0618-satisfiesalphadecayenergylaws}
  \lean{PhyXMiniProblems.ProblemPhyXMini0618.SatisfiesAlphaDecayEnergyLaws}
  \uses{def:physics:phyx-mini-0618:phyxminiproblems-problemphyxmini0618-massinkilograms, def:physics:phyx-mini-0618:phyxminiproblems-problemphyxmini0618-energyinjoules, def:physics:phyx-mini-0618:phyxminiproblems-problemphyxmini0618-speedoflightinmeterspersecond, def:physics:phyx-mini-0618:phyxminiproblems-problemphyxmini0618-conservesnuclearnumbers, def:physics:phyx-mini-0618:phyxminiproblems-problemphyxmini0618-uranium232alphadecaysetup}
  Governing relations for alpha-decay energy: the channel conserves nucleon number and proton number; because the neutral electron counts balance, the mass defect is the parent neutral-atom mass minus the daughter and neutral-helium masses; and the released disintegration energy obeys Q = Delta m c\textasciicircum{}2. No field mentions 5.4 MeV or any answer choice
\end{definition}
\begin{proof}
  This is the declared model definition of Satisfies Alpha Decay Energy Laws.
\end{proof}
\begin{definition}[Answer Choice]
  \label{def:physics:phyx-mini-0618:phyxminiproblems-problemphyxmini0618-answerchoice}
  \lean{PhyXMiniProblems.ProblemPhyXMini0618.AnswerChoice}
  Labels printed beside the four candidate disintegration energies
\end{definition}
\begin{proof}
  This is the declared model definition of Answer Choice.
\end{proof}
\begin{definition}[energy Mega Electron Volts]
  \label{def:physics:phyx-mini-0618:phyxminiproblems-problemphyxmini0618-answerchoice-energymegaelectronvolts}
  \lean{PhyXMiniProblems.ProblemPhyXMini0618.AnswerChoice.energyMegaElectronVolts}
  \uses{def:physics:phyx-mini-0618:phyxminiproblems-problemphyxmini0618-answerchoice}
  Numerical MeV value printed beside each answer label
\end{definition}
\begin{proof}
  This is the declared model definition of energy Mega Electron Volts.
\end{proof}
\begin{definition}[recorded Answer Choice]
  \label{def:physics:phyx-mini-0618:phyxminiproblems-problemphyxmini0618-recordedanswerchoice}
  \lean{PhyXMiniProblems.ProblemPhyXMini0618.recordedAnswerChoice}
  \uses{def:physics:phyx-mini-0618:phyxminiproblems-problemphyxmini0618-answerchoice}
  Dataset metadata records answer C; this is not a theorem premise
\end{definition}
\begin{proof}
  This is the declared model definition of recorded Answer Choice.
\end{proof}
\begin{definition}[Matches Displayed Disintegration Energy]
  \label{def:physics:phyx-mini-0618:phyxminiproblems-problemphyxmini0618-matchesdisplayeddisintegrationenergy}
  \lean{PhyXMiniProblems.ProblemPhyXMini0618.MatchesDisplayedDisintegrationEnergy}
  \uses{def:physics:phyx-mini-0618:phyxminiproblems-problemphyxmini0618-energyinmegaelectronvolts, def:physics:phyx-mini-0618:phyxminiproblems-problemphyxmini0618-answerchoice}
  Agreement with an energy displayed to one decimal place: the calculated value lies within half of the last displayed decimal unit
\end{definition}
\begin{proof}
  This is the declared model definition of Matches Displayed Disintegration Energy.
\end{proof}
% --- Archon declaration topology end ---
% --- Archon physics formalization source end ---
